% Source PDF 6–9, printed vii–x. All four source pages visually checked.
\chapter*{前言}\addcontentsline{toc}{chapter}{前言}
毫无疑问, Alexander Grothendieck 的工作使代数几何发生了革命性的变化. 他引入的许多概念——任意概形、可表函子、相对几何等等——后来都显示出惊人的力量与丰富成果.

1957 至 1962 年间, Grothendieck 在 Séminaire Bourbaki 的一系列报告中勾勒了他的新理论, 并将讲稿汇编成题为“Fondements de la géométrie algébrique”的一卷, 通常简称 FGA. 在 [\textbf{FGA}] 中, 他发展了以下如今已处于绝对核心地位的主题:
\begin{itemize}
\item 下降理论;
\item Hilbert 概形与 Quot 概形;
\item 形式存在定理;
\item Picard 概形.
\end{itemize}
（FGA 还概述了 Grothendieck 将凝聚层的 Serre 对偶推广所得的理论; 这一主题已有相当多的著作详述, 本书不再展开.）

FGA 的许多内容如今已成为常识. 其中一些内容则不那么广为人知, 熟悉其全部范围的几何学家也不多. 然而, 这些理论是代数几何大部分领域中的基本组成部分.

Mudumbai S. Narasimhan 提议在意大利 Trieste 的国际理论物理中心（ICTP）举办一所讲授这些理论的暑期学校. 不过, 这次学校应与 ICTP 的大多数暑期学校有所不同. 后者通常着重介绍当前研究: 讲解重要的新结果, 但证明往往只给梗概或被略去. 这次学校则要传授方法: 证明也必须展开到足够详尽的程度.

2003 年 7 月 7 至 18 日,“Advanced School in Basic Algebraic Geometry”的举办实现了 Narasimhan 的设想. 学术负责人是 ICTP 的 Lother Göttsche、印度 Chennai Mathematical Institute 的 Conjeeveram S. Seshadri, 以及意大利 Università di Bologna 的 Angelo Vistoli. 学校开设了以下课程:
% SOURCE: preface prints Lother, unlike Lothar on title and elsewhere.
\begin{enumerate}[label=(\arabic*)]
\item Angelo Vistoli: Grothendieck 拓扑与下降, 10 学时.
\item Nitin Nitsure: Hilbert 概形与 Quot 概形的构造, 6 学时.
\item Lother Göttsche: Hilbert 概形的局部性质与点的 Hilbert 概形, 4 学时.
\item Luc Illusie: 形式几何中的 Grothendieck 存在定理, 5 学时.
\item Steven L. Kleiman: Picard 概形, 6 学时.
\end{enumerate}
学校主要面向高年级研究生, 也面向刚开始研究的学者; 两类学员都热情参与. ICTP 的行政组织十分专业. 每个人都有一段难忘的经历.

% Source PDF 7, printed viii.
本书分为五个部分, 是学校所发讲义的扩充、修订版本. 本书无意取代 [\textbf{FGA}]; 事实上, 大师本人的文字永远无可替代, 阅读 Grothendieck 总能给人启迪. 本书所做的是补全 Grothendieck 的纲要. 此外, 只要较新的思想有助于理解, 本书就予以引入, 并指出与后续发展的联系. 例如, 本书采用 Grothendieck 后来引入的 Grothendieck 拓扑语言来叙述下降理论. Hilbert 概形和 Picard 概形的有限性是困难的基本结果; 其证明不采用 Chow 坐标, 而采用 Castelnuovo--Mumford 正则性, 后者如今已是代数几何与交换代数中的主要工具.

本书无意提供迅速、轻松的入门. 相反, 书中给出的是要求较高的详细论述. 其回报是对材料深入得多的理解. 主要先修要求是透彻熟悉教材 [\textbf{Har77}] 所发展的基本概形论.

本书内容简述如下. 五个部分的引言中各有更详尽的概述.

第一部分由 Vistoli 撰写, 对下降理论作了相当完整的论述. 它既解释抽象方面——纤维范畴与叠——也解释最重要的具体情形——拟凝聚层与概形的下降. 第一部分包括第 1 至 4 章.

第 1 章回顾范畴论与代数几何的一些基本概念. 第 2 章介绍可表函子、Grothendieck 拓扑和层; 这些概念广为人知, 已有几种很好的论述, 但本书的处理或许对初学者更有吸引力, 也可作为随后较高深理论的热身.

第 3 章专门讨论一个基本概念——纤维范畴, 这是 Grothendieck 在 [\textbf{SGA1}] 中引入的. 主要例子是概形范畴上的拟凝聚层范畴. 纤维范畴为讨论下降理论提供了合适的抽象框架. 尽管初等应用可能不需要一般理论, 但深入理解与高阶应用却离不开它.

第 4 章讨论叠, 即下降理论成立的纤维范畴. 本章完整讨论定义下降数据的不同方式, 并证明第一部分的主要结果: 拟凝聚层构成一个叠.

第二部分由 Nitsure 撰写, 沿着 Grothendieck 的 Bourbaki 报告 [\textbf{FGA}, 221], 介绍 Hilbert 概形和 Quot 概形的构造, 同时介绍 David Mumford 以及 Allen Altman 与 Kleiman 的后续发展. 第二部分包括第 5 章.

具体而言, 给定概形 $X$, Grothendieck 解决了如下基本问题: 构造另一个概形 $\Hilb_X$, 称为 $X$ 的 \emph{Hilbert 概形}, 以适当的泛性质参数化 $X$ 的所有可能的闭子概形. 更一般地, 给定 $X$ 上的凝聚层 $E$, 他构造了一个概形 $\Quot_{E/X}$, 称为 $E$ 的 \emph{Quot 概形}, 同样以适当的泛性质参数化 $E$ 的所有可能的凝聚商层. 在相对情形下, 当 $X$ 在适当的底上射影时, 可以进行这些构造. 构造的关键在于运用几个基本工具, 包括忠实平坦下降、平坦化分层、半连续性复形, 以及 Castelnuovo--Mumford 正则性.

% Source PDF 8, printed ix.
第三部分由 Barbara Fantechi 与 Göttsche 合写, 包括第 6、7 两章.

第 6 章介绍（无穷小）形变函子的概念, 并给出几个例子. 本章还为这样的函子定义切空间—障碍理论, 解释这一理论如何给出模空间维数的估计. 对某些情形给出完整论述, 对另外几个第 7 章不需要的情形则只作概述.

第 7 章研究光滑拟射影簇上点的 Hilbert 概形, 它参数化固定长度的有限子概形. 本章构造 Hilbert--Chow 态射, 将一个子概形映到带重数的支集, 从而把此 Hilbert 概形映到对称幂. 对曲面, 这个态射给出奇点解消. 最后, 本章计算 Hilbert 概形的 Betti 数, 并概述 Heisenberg 代数在上同调上的作用.

第四部分由 Illusie 撰写, 重访 Grothendieck 的 Bourbaki 报告 [\textbf{FGA}, 182]. 在该报告中, 他提出代数几何与形式几何之间一个基本的“GAGA”型比较定理, 并概述它对基本群理论和无穷小形变理论的某些应用. 不久之后, 详细论述见于 [\textbf{EGAIII1}]、[\textbf{EGAIII2}] 和 [\textbf{SGA1}]. 第四部分包括第 8 章.

第 8 章先回顾局部 Noether 形式概形的基本事实, 再解释主要比较定理证明中的关键点, 并概述一些推论, 包括 Zariski 连通性定理与主定理, 以及 Grothendieck 关于形式概形代数化的判据. 随后介绍 Grothendieck 对基本群以及向量丛和光滑概形提升问题的应用, 特别是曲线与 Abel 簇的情形. 本章最后讨论 Serre 的著名例子: 正特征的簇不能提升到特征零.

第五部分由 Kleiman 撰写, 详细发展 Grothendieck 在两篇 Bourbaki 报告 [\textbf{FGA}, 232, 236] 及其评论 [\textbf{FGA}, 第 C-07--C-011 页] 中概述的 Picard 概形理论的大部分内容. 此外, 第五部分还以散见各处的一系列注, 简要回顾 Grothendieck 与他人发展的其余大部分理论. 第五部分包括第 9 章.

第 9 章以一篇详尽的历史引言开头, 通过追溯有关思想的发展来说明 Grothendieck 研究 Picard 概形的动机. 这段历史引人入胜, 本身或许也值得独立阅读.

随后, 第 9 章讨论四种常见的相对 Picard 函子, 它们依次更有可能成为 Picard 概形的点函子. 接着讨论两个重要的预备概念: 相对有效（Cartier）除子和线性等价. 然后证明 Grothendieck 关于 Picard 概形的主定理: 对任意几何纤维均为整概形的射影平坦概形, Picard 概形存在.

接下来, 第 9 章研究 Picard 概形各纤维单位元所在的连通分支之并. 然后证明两个更深的有限性定理: 第一个涉及某个倍数落入此并集的那些点; 第二个涉及表示具有给定 Hilbert 多项式的可逆层的那些点. 第 9 章以两个附录结束: 一个给出全部习题的详细解答; 另一个以初等方式介绍基本的除子相交理论——这套理论在两个有限性定理的证明中被自由使用.
% Last paragraph spans source PDF 8–9.
